%% kapa: exam=1 round=24 year=2013 period=2 form=A subject=appraisal_law subject_ko=부동산관계법규 seq=24 number=64
%% source: Q-Net 감정평가사 기출문제 (cst00309) artlSeq=5073802 file=extracted/1008/5073802/2013년 제24회 감정평가사 1차시험 기출문제/2교시(A형).hwp
%% conversion: kordoc 4.24.51 -> tools/gen1.py (auto); review=proofread-sonnet; flags=none
\begin{question}{64}
부동산 가격공시 및 감정평가에 관한 법령상 표준지공시지가에 관한 설명으로 옳지 않은 것은?
\begin{choices}
  \item 표준지의 도로ㆍ교통상황, 지세(地勢)는 표준지공시지가의 공시사항에 포함된다.
  \item 지상권이 설정되어 있는 표준지의 적정가격을 평가하는 경우에는 당해 지상권으로 토지의 사용ㆍ수익이 제한되는 가치를 반영하여 평가하여야 한다.
  \item 표준지의 단위면적당 가격에서 단위면적은 1제곱미터로 한다.
  \item 표준지공시지가는「도시개발법」에 의한 도시개발사업을 위한 환지ㆍ체비지의 매각 또는 환지신청에 적용된다.
  \item 표준지공시지가결정이 위법한 경우에는 항고소송으로 그 위법 여부를 다툴 수 있다.
\end{choices}
\end{question}
\answer{2}
